
\documentclass[12pt,a4paper]{article}

% --------------------------------------------------
% Codificación e idioma
% --------------------------------------------------
\usepackage[utf8]{inputenc}
\usepackage[T1]{fontenc}
\usepackage[spanish,es-noquoting]{babel}
\usepackage{lmodern}
\usepackage{microtype}

% --------------------------------------------------
% Diseño de página
% --------------------------------------------------
\usepackage[
    a4paper,
    left=3cm,
    right=3cm,
    top=3cm,
    bottom=3cm
]{geometry}

\usepackage{setspace}
\setstretch{1.25}

\usepackage{parskip}
\usepackage{csquotes}
\usepackage{hyperref}

\hypersetup{
    colorlinks=true,
    urlcolor=blue,
    linkcolor=black,
    pdftitle={Del poder y el control en la desigualdad.
              De un gavilán y un aeda},
    pdfauthor={Vicente Domínguez Arca}
}

% --------------------------------------------------
% Datos del documento
% --------------------------------------------------
\title{
    \textbf{Del poder y el control en la desigualdad}\\[0.4em]
    \large De un gavilán y un aeda
}

\author{
    Vicente Domínguez Arca
}

\date{}

% --------------------------------------------------
% Documento
% --------------------------------------------------
\begin{document}

\maketitle

\vspace{1em}

Un lector atento seguramente se habrá percatado de que
tras este proyecto doctoral existe una posición intelectual
muy clara por parte de su autor. No pretendo ocultarla;
nunca lo he pretendido. Pero entiendo que, en este
desarrollo, es necesario afanarse en la neutralización
del prejuicio, no con el fin de hacer concesiones, sino
con la firme intención de lograr que las conclusiones
y los argumentos descansen sobre cimientos sólidos.

En esta nota reflexiva, particularmente personal,
he preferido trasladarme a la literatura clásica
y rescatar unas palabras de Hesíodo en su obra
\textit{Los trabajos y los días}, concretamente
en el libro I:

\begin{quote}
\itshape
Un gavilán habló así a un ruiseñor sonoro al que
había cogido en sus garras y se lo llevaba por
las altas nubes. El ruiseñor, desgarrado por las
curvas uñas, gemía; pero el gavilán le dijo estas
palabras imperiosas:

\medskip

Desdichado, ¿por qué gimes? Ciertamente, eres
presa de uno más fuerte que tú. Irás adonde
yo te conduzca, aunque seas un aeda.
Te comeré, si me place, o te soltaré.
¡Malhaya quien quiera luchar contra otro
más poderoso que él! Será privado de la
victoria y abrumado de vergüenza y de dolores.

\medskip

Así habló el rápido gavilán de anchas alas.
\end{quote}

La parábola siempre me ha parecido maravillosa en cada una
de sus palabras. Es extraordinariamente densa:
todo parece encerrar un significado mucho más
profundo de lo que muestra a primera vista y,
como sucede con buena parte de la literatura
clásica, posee esa inexplicable esencia de
atemporalidad.

Siempre me ha llevado a pensar en lo que considero
una de las mayores paradojas del poder.

El gavilán de anchas alas se siente tan poderoso
sobre el ruiseñor que parece encontrar en sus
garras y en sus altos vuelos entre las nubes
la expresión indiscutible de su dominio.
Sin embargo, siempre he supuesto que también
es plenamente consciente de las necesidades
de su hambriento estómago.

Y me pregunto: ¿qué sería del gavilán sin
el ruiseñor?

¿De quién es el poder, entonces?

¿Existe tal cosa como el poder?

Quizás sea una reflexión demasiado abierta,
incluso demasiado ambiciosa. Pero si en la
jerarquía de la desigualdad, el control y el
poder introducimos el factor de la necesidad,
en todos y cada uno de los escalones de esa
pavorosa verticalidad, ¿qué ocurre?

¿Qué se nos muestra cuando observamos desde la necesidad?


\end{document}

